% Reproducible case-level audit; the protocol here uses greedy B=1.
\subsection{Case-level Success, Failure, and Structural Diversity}
\label{appendix:case_diagnostics}
\paragraph{Diagnostic protocol.}
We analyze seven deliberately selected C-MuMOInstruct test cases
(\texttt{instr\_setting=seen}). For Mistral and Llama, we compare
GeLLM$^4$O-C with the GRPO and GDPO post-trained checkpoints;
for MolT5 we compare its SFT, GRPO, and GDPO checkpoints.
These examples use \emph{greedy decoding with a single candidate},
not the beam-20 candidate-selection protocol used for the main
benchmark; therefore, individual outcomes must not be used as
estimates of the main-table success rates. Generation seed 42 was
recorded for the Mistral/Llama generations, whereas the archived
MolT5 outputs have no recorded seed. SOR/SSOR (denoted 1/0 below)
and per-property pass/fail decisions were recalculated with the
project's test property calculators and source-dependent
$\Delta_p$ rules. ADMET predictions are not experimental measurements.
Similarity is Morgan radius-2/2048-bit Tanimoto and is descriptive:
the case-level evaluation imposes no explicit similarity cutoff.
This intentionally contrastive sample is not statistically
representative and does not identify the causal effect of an
optimization objective on a particular edit.

\begin{table*}[t]
\centering\footnotesize
\setlength{\tabcolsep}{4pt}\renewcommand{\arraystretch}{1.14}
\begin{tabularx}{\textwidth}{@{}llcccX@{}}
\toprule
\textbf{Case ID} & \textbf{Backbone} & \textbf{Base}
& \textbf{GRPO} & \textbf{GDPO} & \textbf{Illustrated outcome}\\
\midrule
BCMQ \#19 & Mistral & 0/0 & \textbf{1/1} & 0/0 &
GRPO succeeds; GDPO fails CARC and MUTAG preservation.\\
ACEP \#127 & Mistral & 0/0 & 0/0 & \textbf{1/1} &
Only GDPO achieves the required CARC decrease.\\
ABMP \#343 & Mistral & 1/0 & 1/0 & \textbf{1/1} &
Both RL outputs satisfy SOR but differ in molecule and SSOR.\\
BCMQ \#3 & Llama & 0/0 & \textbf{1/1} & 1/0 &
Both satisfy SOR; GRPO obtains higher QED and SSOR.\\
BCMQ \#333 & Llama & 0/0 & \textbf{1/1} & \textbf{1/1} &
Both satisfy strict success but produce different molecules.\\
ABMP \#18 & Mistral & 0/0 & 0/0 & 0/0 &
GRPO fails MUTAG improvement; GDPO fails BBBP preservation.\\
ABMP \#417 & MolT5 & 0/0 & 0/0 & \textbf{1/1} &
GDPO satisfies property criteria but similarity is only 0.0476.\\
\bottomrule
\end{tabularx}
\caption{Seven contrasting, \emph{single-candidate} examples.
Entries are SOR/SSOR (1 = pass, 0 = fail).
Base means GeLLM$^4$O-C for Mistral/Llama and SFT for MolT5.
Within each row all methods use the same source and task.}
\label{tab:case_diagnostic_overview}
\end{table*}

\paragraph{Success and failure are often determined by one bottleneck.}
In \textbf{BCMQ \#19 (Mistral)}, GRPO reaches 1/1 whereas
GDPO remains at 0/0. With source CARC=0.1345 (a property to
preserve), GDPO predicts 0.4112, exceeding the permitted
0.2 degradation; GRPO instead predicts 0.0186.
GDPO also raises the MUTAG prediction from 0.1462 to 0.2612,
slightly beyond its permitted 0.1 degradation, while GRPO
predicts 0.0768.
The reverse occurs for \textbf{ACEP \#127 (Mistral)}:
the source CARC=0.3452 requires a decrease of at least 0.2,
i.e., a prediction no greater than 0.1452.
GeLLM$^4$O-C, GRPO, and GDPO predict CARC values of
0.2071, 0.1636, and 0.0615, respectively; only GDPO
clears the improvement margin, attaining 1/1.
The GDPO candidate also differs chemically, replacing the
amide-containing local substituent in the GRPO output
with a carbonyl-linked cyclopropyl substituent.
This describes an observed edit and predicted-score difference,
not a validated medicinal-chemistry mechanism.

\paragraph{Both GRPO and GDPO can succeed on different molecules.}
For \textbf{ABMP \#343 (Mistral)}, both outputs satisfy SOR
but differ in chlorophenyl substitution. GRPO obtains
BBBP=0.8654 and similarity=0.6269 (SOR/SSOR=1/0);
GDPO obtains BBBP=0.9127 and similarity=0.6769
(SOR/SSOR=1/1). Both clear the BBBP source-relative
improvement margin from 0.7017, but only GDPO reaches
the 0.9 strict target. For \textbf{BCMQ \#3 (Llama)},
GRPO and GDPO again generate different molecules
with SOR=1: GRPO achieves QED=0.9185 and SSOR=1,
versus QED=0.7879 and SSOR=0 for GDPO.
Finally, in \textbf{BCMQ \#333 (Llama)} both methods
reach SOR=SSOR=1 using two \emph{distinct} products:
the GRPO product has similarity 0.4898, whereas the
GDPO product has similarity 0.3529.
Thus, even identical binary success outcomes do not imply
identical edits or equally strong structural conservation.
Table~\ref{tab:case_smiles_diagnostic} gives their exact SMILES.

\paragraph{Two failures, two different violated objectives.}
In \textbf{ABMP \#18 (Mistral)}, both RL candidates are valid
SMILES and both fail SOR, but for different reasons.
GRPO decreases MUTAG from 0.3387 to 0.2728, a gain
of only 0.0659 compared with the required 0.1 reduction.
GDPO reaches MUTAG=0.0504, yet decreases BBBP from 0.9099
to 0.6743, exceeding the permitted 0.1 degradation.
The failure cannot be understood from a single scalar
reward or overall SOR alone.

\paragraph{Property success can mask drastic structural change.}
In \textbf{ABMP \#417 (MolT5)}, GDPO attains SOR=SSOR=1
with predicted AMP=0.9988, BBBP=0.9540, MUTAG=0.0005,
and PlogP=4.0041. The product is a heavily alkyl-substituted
benzene, radically different from the heteroatom-rich
multi-ring source (Tanimoto=0.0476); SFT and GRPO both
fail the MUTAG constraint. This case reveals structural
drift permitted by an evaluation based on property criteria
without a hard scaffold-preservation constraint.
It is \emph{not} evidence of experimental molecular quality.
Reporting similarity alongside SOR and SSOR is therefore
essential, even when all scored properties pass.

\begin{table*}[t]
\centering\small
\setlength{\tabcolsep}{5pt}\renewcommand{\arraystretch}{1.10}
\begin{tabular}{@{}llrrrrl@{}}
\toprule
\textbf{Case} & \textbf{Property} & \textbf{Source}
& \textbf{Base} & \textbf{GRPO} & \textbf{GDPO}
& \textbf{Key failure / contrast}\\
\midrule
BCMQ \#19 & CARC $\downarrow$ & .1345 & .5348 & .0186 & .4112 & GDPO preservation\\
BCMQ \#19 & MUTAG $\downarrow$ & .1462 & .0881 & .0768 & .2612 & GDPO preservation\\
ACEP \#127 & CARC $\downarrow$ & .3452 & .2071 & .1636 & .0615 & GRPO below improvement margin\\
ABMP \#343 & BBBP $\uparrow$ & .7017 & .8783 & .8654 & .9127 & Only GDPO at strict target\\
BCMQ \#3 & QED $\uparrow$ & .8080 & .8080 & .9185 & .7879 & Different strict success\\
ABMP \#18 & MUTAG $\downarrow$ & .3387 & .2805 & .2728 & .0504 & GRPO insufficient improvement\\
ABMP \#18 & BBBP $\uparrow$ & .9099 & .9114 & .9931 & .6743 & GDPO preservation loss\\
ABMP \#417 & MUTAG $\downarrow$ & .3599 & .6472 & .3814 & .0005 & GDPO success, severe structural drift\\
\bottomrule
\end{tabular}
\caption{Source and output predicted property values for
diagnostic cases. Up/down indicate preferred property direction;
whether a property is improved or preserved is conditioned on
the source value under the benchmark's $\Delta_p$ rule.}
\label{tab:case_diagnostic_properties}
\end{table*}

\begin{table*}[t]
\centering\footnotesize
\setlength{\tabcolsep}{4pt}\renewcommand{\arraystretch}{1.15}
\begin{tabularx}{\textwidth}{@{}llcX@{}}
\toprule
\textbf{Case} & \textbf{Candidate} & \textbf{SOR/SSOR} &
\textbf{Full SMILES}\\
\midrule
\multirow{4}{*}{ACEP \#127 / Mistral}
& Source & --- & \nolinkurl{Cc1ccc(Cn2nc(C)c(C(=O)NCc3cc(Cl)c4c(c3)OCCCO4)c2Cl)cc1}\\
& GeLLM$^4$O-C & 0/0 & \nolinkurl{Cc1ccc(Cn2nc(C)c(C(=O)NCC3CC3)c2Cl)cc1}\\
& GRPO & 0/0 & \nolinkurl{Cc1ccc(Cn2nc(C)c(C(=O)NC3CC3)c2Cl)cc1}\\
& GDPO & 1/1 & \nolinkurl{Cc1ccc(Cn2nc(C)c(C(=O)C3CC3)c2Cl)cc1}\\
\midrule
\multirow{4}{*}{BCMQ \#3 / Llama}
& Source & --- & \nolinkurl{Cc1cccnc1NC(=O)CC(C)(C)C}\\
& GeLLM$^4$O-C & 0/0 & \nolinkurl{Cc1cccnc1NC(=O)CC(C)(C)C}\\
& GRPO & 1/1 & \nolinkurl{Cc1cccnc1NC(=O)CC(C)(C)N1CCCCC1}\\
& GDPO & 1/0 & \nolinkurl{Cc1cccnc1NC(=O)N1CCCCC1}\\
\midrule
\multirow{4}{*}{BCMQ \#333 / Llama}
& Source & --- & \nolinkurl{c1ccc([C@H](CCNCc2ccco2)c2ccco2)cc1}\\
& GeLLM$^4$O-C & 0/0 & \emph{Invalid generated SMILES}\\
& GRPO & 1/1 & \nolinkurl{c1ccc([C@H](CCNC2CCN(C)CC2)c2ccco2)cc1}\\
& GDPO & 1/1 & \nolinkurl{CC(C)CNC1CC([C@H](c2ccccc2)c2ccco2)C1}\\
\bottomrule
\end{tabularx}
\caption{Audited source and generated SMILES for three
within-backbone examples, showing that GRPO and GDPO can
produce different molecules despite the same SOR.
The baseline SMILES in BCMQ \#333 was invalid and is omitted.}
\label{tab:case_smiles_diagnostic}
\end{table*}

\paragraph{Summary.}
These examples show that a single missed property can determine
binary success, that equal SOR can conceal different molecules
and strict-success outcomes, and that property success alone
does not guarantee structural preservation. Neither optimizer
dominates every selected case; these qualitative observations
complement, rather than replace, the full benchmark statistics.
